% ── SECTION 6b: THE RCT IN PHILOSOPHICAL LANGUAGE (~900 words) ───────────────

\section{The Relational Coherence Theorem: Formalizing the Convergence}
\label{sec:rct-phil}

The three preceding sections have made a philosophical argument.
Sections~\ref{sec:substance}--\ref{sec:whitehead} established that quantum
mechanics requires process ontology and that Whitehead's framework names
its structure.
Sections~\ref{sec:bahai}--\ref{sec:convergences} demonstrated that the Bahá'í
writings constitute an independent prior instantiation of the same structure
across six specific convergences.
Section~\ref{sec:semantic} showed that each tradition's technical vocabulary
names the same underlying phenomenon.

What converts this philosophical argument into a scientific contribution
with testable consequences is the \emph{Relational Coherence Theorem} (\RCT):
a formal mathematical framework, developed in the companion physics
paper~\citep{Adams2026Physics}, that translates the central claim of
process ontology into the language of quantum mechanics and derives
quantitatively distinct predictions from it.
This section presents the \RCT\ in non-mathematical terms sufficient for
the philosophical argument.

\subsection{What the RCT Does}

The central claim of process ontology --- and of the three-tradition
convergence documented in this paper --- is that \emph{relational affinity
is the condition of existence}.
Things exist not as independent substances possessing intrinsic properties,
but as patterns of relational coherence: the ongoing relational bonds
that constitute a thing are what make it a thing.
When those bonds dissolve, the thing dissolves.

The \RCT\ translates this claim into formal quantum mechanics through
a single quantitative object: the \emph{Affinity Tensor} $\alpha(S,E)$.
For any quantum system $S$ in relation to its environment $E$, the
Affinity Tensor measures the degree to which $S$ retains relational
coherence with $E$ --- the magnitude of the quantum bonds that constitute
$S$'s participation in relational existence.
When $\alpha = 1$, $S$ is maximally coherent: fully participating in
the relational structure that constitutes quantum existence.
When $\alpha = 0$, $S$ is fully decohered: classically separate,
its relational bonds dissolved into the environment.
The range $[0,1]$ is not a metaphor.
It is a precisely defined, experimentally accessible, continuously
varying measure of what process philosophy calls the degree of a
thing's relational being.

\subsection{Four Propositions}

The \RCT\ comprises four propositions, each of which makes the
philosophical claim mathematically precise.

\paragraph{Property indefiniteness.}
A system with nonzero Affinity Tensor ($\alpha > 0$) does not possess
a definite property in the classical sense.
This is not a statement about measurement uncertainty.
It is a statement about ontological structure: the system is still in
the process of becoming, still in the superposition of relational
possibilities that Whitehead would call its potential phase.
Definite properties belong to decohered states; relational coherence
and classical definiteness are mutually exclusive, exactly as process
ontology predicts.

\paragraph{Relational determination.}
Definite properties emerge through relational interaction.
Under environmental coupling, the Affinity Tensor decays continuously
from its initial value toward zero.
As it decays, the system's relational potential is actualized into
definite outcomes through its interaction with the environment.
This is the mathematical description of Whiteheadian concrescence:
the transition from potential to actuality through relational process.

\paragraph{Coherence--affinity correspondence.}
The Affinity Tensor ranges continuously over $[0,1]$, with $\alpha = 1$
at maximum relational coherence and $\alpha = 0$ at full decoherence.
This means that relational affinity is not binary --- not simply present
or absent --- but a continuously graded quantity.
Existence, in the process ontological sense, admits of degrees.
The Bahá'í passage that defines nonexistence as the \emph{absence} of
affinity and attraction (Section~\ref{sec:bahai}) is, on this reading,
a statement about the $\alpha = 0$ limit: the point at which relational
coherence has been fully extinguished and the pattern of relational
events that constituted the thing has dissolved.

\paragraph{Conservation of relational potential.}
When a quantum system is not interacting with its environment --- when
it evolves in isolation --- its total relational potential is conserved.
What changes is its distribution: relational coherence that was locally
concentrated in the system is redistributed into correlations between
the system and its environment.
Relational being is not created or destroyed; it is transformed.
This is the formal expression of what the Bahá'í writings call the
creative foundation: the relational ground of existence is not consumed
by the process of becoming but redistributed through it, available for
the next occasion of relational actualization.

\subsection{The Relational Affinity Gate}

The \RCT\ has a direct application in quantum computing.
If relational coupling is a continuously variable resource rather than
a fixed quantity, then a two-qubit gate should expose a continuous
coupling parameter that can be tuned to the actual coherence budget of
a physical quantum processor.
The \emph{Relational Affinity Gate} ($\URA$), derived from the \RCT,
does exactly this: it is a continuously parameterized two-qubit unitary
that trades a fraction of its entangling capability for a proportionally
shorter operation time under noise.
On hardware where gate duration scales with coupling strength, $\URA$
at sub-maximal coupling achieves measurably higher process fidelity
than fixed entangling gates such as \textsc{cnot} or \textsc{iswap}.
This is not merely an engineering optimization.
It is a prediction derived from the process ontological claim that
relational affinity is a continuous, conservatively redistributed
resource --- a claim that the \RCT\ formalizes and that the $\URA$
gate tests experimentally.

\subsection{The RCT as the Formal Expression of the Convergence}

The significance of the \RCT\ for the argument of this paper is this:
it demonstrates that the convergence documented in the preceding sections
is not merely philosophical.
The claim that existence is constituted by relational affinity --- made
by Whitehead through philosophical analysis, by the Bahá'í writings
through spiritual revelation, and by quantum mechanics through a century
of experimentation --- is sufficiently precise to generate a formal
mathematical framework, a novel quantitative object (the Affinity
Tensor), a new two-qubit gate, and testable predictions about gate
fidelity under realistic noise.

The convergence is not a set of interesting similarities.
It is the identification of a real feature of reality --- the relational
constitution of existence --- that is precise enough to be formalized,
computed, and tested.
The \RCT\ is that formalization.
